\documentclass[modern,large,numbib]{oup-authoring-template}

\usepackage{amsmath,amssymb}
\usepackage{graphicx}
\usepackage{booktabs}
\usepackage{multirow}
\usepackage{algorithm,algpseudocode}
\usepackage{hyperref}

\journalname{Bioinformatics}
\doi{XX.XXXX/XXXXXXXXXX}
\copyrightyear{2026}

\begin{document}

\title{BADGE: Bi-iterative Adaptive Drug-disease Graph Enhancement for Drug Repositioning}

\author{Yankang Y.}

\maketitle

\begin{abstract}
\textbf{Motivation:} Existing matrix completion methods for drug repositioning
treat similarity kernels as fixed exogenous inputs, computed once from sparse
known associations and then held constant throughout optimization. However, the
completion matrix and the similarity kernel are mutually dependent: a better
completion yields a more informative kernel, which in turn enables better
completion. This interdependence motivates a joint estimation framework where
both the completion and the kernel are optimization variables.

\textbf{Results:} We propose BADGE (Bi-iterative Adaptive Drug-disease Graph
Enhancement), which formalizes drug-disease matrix completion as a joint
matrix-kernel estimation problem and solves it via alternating minimization.
Starting from GIP kernels derived from sparse associations (iteration 1,
equivalent to GF-BNNR), BADGE recomputes the kernel from the filtered completion
and performs a second round of completion (iteration 2). The alternating scheme
converges at $N=2$ iterations: the GIP kernel stabilizes after one refinement
because the filtered completion is already near the fixed point of the kernel
update map. BADGE nests prior methods as special cases---BNNR \texttt{Demo2} and
GF-BNNR correspond to one iteration without and with the graph filter,
respectively. Across three benchmark datasets spanning three orders of magnitude
in association density (0.015\%--1.04\%), BADGE achieves the highest AUPR on all
three, improving over the BNNR baseline by $+5.6\%$ (Fdataset), $+46.1\%$
(Cdataset), and $+25.1\%$ (DNdataset). Ablation confirms that both the
endogenous GIP kernel and the graph filter are essential---removing either causes
severe degradation on ultra-sparse data.

\textbf{Availability:} Code available at
\url{https://github.com/yankangyk/bnnr-innovation}.
\end{abstract}

% ===========================================================================
\section{Introduction}
% ===========================================================================

Drug repositioning---identifying new therapeutic indications for approved
drugs---offers substantial advantages over de novo drug discovery by leveraging
existing safety and pharmacokinetic data \cite{chong2007,paul2010}.
Computational approaches formulate this as a matrix completion problem: given a
partially observed drug-disease association matrix and drug-drug and
disease-disease similarity matrices, predict the unknown entries
\cite{yang2019,luo2018,liu2016,zheng2013}.

Bounded Nuclear Norm Regularization (BNNR; \cite{yang2019}) is a widely adopted
matrix completion framework that integrates drug and disease similarities into a
heterogeneous network adjacency matrix and recovers missing entries via nuclear
norm minimization with a range constraint, solved by the Alternating Direction
Method of Multipliers (ADMM). BNNR tolerates noise in similarity computations
and handles cold-start scenarios effectively.

BNNR integrates similarity matrices through an augmented heterogeneous adjacency
matrix (Eq.~\ref{eq:T}). While the original paper uses raw chemical and clinical
similarities, the accompanying code release (\texttt{Demo2.m} in
\cite{yang2019}) also demonstrates fusing GIP kernels with raw similarities:
$\tilde{\mathbf{S}} = w_{gip} \mathbf{G} + (1 - w_{gip}) \mathbf{S}_{raw}$,
showing that association-derived structure improves completion accuracy.
However, in this standard approach the GIP kernel is computed \textit{once} from
the sparse observation matrix $\mathbf{W}_{dr}$ and then fixed as an exogenous
input throughout the optimization. The kernel quality is therefore bounded by
the sparsity of the observed data.

We previously developed two extensions to incorporate manifold geometry more
explicitly: \textbf{GBNNR}, which injects graph Laplacian regularization into
the ADMM iterations, and \textbf{GF-BNNR}, which applies a post-hoc
bi-directional graph low-pass filter. Both methods substantially improve over
BNNR (Table~\ref{tab:main}), but share the same structural limitation as
\texttt{Demo2}: the similarity kernel is treated as a fixed exogenous input.

Here we recognize a deeper issue: the similarity kernel $\mathbf{S}$ and the
completion matrix $\mathbf{M}$ are mutually dependent. A better completion
enables a more accurate GIP kernel; a better kernel enables better completion.
This interdependence suggests a \textbf{joint estimation} problem rather than
the conventional fixed-kernel formulation. We formalize this as:
%
\begin{equation}
\min_{\mathbf{M}, \mathbf{S}_d, \mathbf{S}_t}
\|\mathbf{M}\|_* + \frac{\alpha}{2}\|P_\Omega(\mathbf{M} - \mathbf{Y})\|_F^2
+ \alpha_f \cdot
\big(\text{tr}(\mathbf{M}^\top \mathbf{L}_t \mathbf{M})
+ \text{tr}(\mathbf{M} \mathbf{L}_d \mathbf{M}^\top)\big)
\label{eq:joint_obj}
\end{equation}
%
\begin{equation}
\text{s.t.}\quad
\mathbf{S}_d = w_{gip} \cdot \text{GIP}_{\text{drug}}(\mathbf{M})
+ (1 - w_{gip}) \cdot \mathbf{S}_{rr},
\quad
\mathbf{S}_t = w_{gip} \cdot \text{GIP}_{\text{dis}}(\mathbf{M})
+ (1 - w_{gip}) \cdot \mathbf{S}_{dd}
\label{eq:kernel_constraint}
\end{equation}
%
where $\mathbf{L}_d, \mathbf{L}_t$ are the normalized graph Laplacians of
$\mathbf{S}_d, \mathbf{S}_t$, and $\alpha_f \geq 0$ controls the graph filter
strength. Critically, Eq.~\ref{eq:kernel_constraint} makes the kernel an
\textit{endogenous} variable that depends on the current completion $\mathbf{M}$,
rather than an exogenous input computed once from sparse observations.

Problem~(\ref{eq:joint_obj})--(\ref{eq:kernel_constraint}) is biconvex: for
fixed $\mathbf{S}_d, \mathbf{S}_t$, the $\mathbf{M}$-subproblem is convex
(nuclear norm + quadratic + graph regularization); for fixed $\mathbf{M}$, the
kernel update has the closed form~(\ref{eq:kernel_constraint}). We solve it via
alternating minimization (Algorithm~\ref{alg:badge}), which converges to a
stationary point under standard conditions \cite{tseng2001}.

We term this framework BADGE (Bi-iterative Adaptive Drug-disease Graph
Enhancement). The standard one-shot GIP approach (BNNR \texttt{Demo2}) and
GF-BNNR emerge as special cases: one iteration with $\alpha_f=0$ and
$\alpha_f>0$, respectively. With $N=2$ iterations, BADGE performs one round of
kernel refinement using the filtered completion from iteration 1, after which
the kernel stabilizes and further iterations yield negligible change.

The main contributions of this work are:
\begin{itemize}
    \item \textbf{Joint matrix-kernel estimation.} We formalize drug-disease
    matrix completion as a joint estimation problem over the completion
    $\mathbf{M}$ and similarity kernel $\mathbf{S}$, where the kernel is an
    endogenous function of the completion rather than a fixed exogenous input.
    This reformulates the role of GIP from a preprocessing step to an integral
    part of the optimization.
    \item \textbf{BADGE algorithm.} An alternating minimization procedure that
    solves the joint problem, achieving the highest AUPR across all three
    benchmark datasets. Relative to the BNNR baseline \cite{yang2019}, BADGE
    improves AUPR by $+5.6\%$ (Fdataset), $+46.1\%$ (Cdataset), and $+25.1\%$
    (DNdataset).
    \item \textbf{Convergence analysis.} We prove that the alternating
    minimization converges to a stationary point of the biconvex objective, and
    characterize why $N=2$ iterations suffice in practice: the GIP kernel
    stabilizes after one round of refinement because the filtered completion is
    already near the fixed point of the kernel update operator.
    \item \textbf{Progressive development} from BNNR (raw similarities) through
    \texttt{Demo2} (one-shot GIP), GBNNR (inside graph regularization), GF-BNNR
    (outside graph filtering), to BADGE (joint estimation with iterative kernel
    refinement), with each stage providing cumulative improvement.
    \item \textbf{Ablation and sensitivity analyses} establishing that GIP
    fusion and graph filtering are individually essential, and that
    $w_{gip}=0.3$ is a robust default validated by systematic sweep across all
    three datasets.
\end{itemize}

% ===========================================================================
\section{Materials and Methods}
% ===========================================================================

\subsection{Datasets}

We evaluate on three benchmark datasets widely used in drug repositioning
research \cite{gottlieb2011,yang2019}:

\begin{table}[ht]
\centering
\caption{Benchmark dataset statistics.}
\label{tab:datasets}
\begin{tabular}{lrrrr}
\toprule
Dataset & Drugs & Diseases & Known Associations & Density \\
\midrule
Fdataset  & 593  & 313  & 1,933 & 1.04\% \\
Cdataset  & 663  & 409  & 2,532 & 0.93\% \\
DNdataset & 1,490 & 4,516 & 1,008 & 0.015\% \\
\bottomrule
\end{tabular}
\end{table}

Drug-drug similarities are computed as Tanimoto scores between chemical
fingerprints derived from SMILES representations (DrugBank), using the Chemistry
Development Kit \cite{steinbeck2003}. Disease-disease similarities are obtained
from MimMiner \cite{vandriel2006}, which measures overlap of MeSH terms in OMIM
disease descriptions. All similarity values lie in $[0, 1]$.

\subsection{Preliminaries}

\subsubsection{BNNR \cite{yang2019}}

Let $\mathcal{R} = \{r_1, \ldots, r_m\}$ be $m$ drugs and
$\mathcal{D} = \{d_1, \ldots, d_n\}$ be $n$ diseases. The drug-drug similarity
$\mathbf{S}_{rr} \in [0,1]^{m \times m}$, disease-disease similarity
$\mathbf{S}_{dd} \in [0,1]^{n \times n}$, and drug-disease association
$\mathbf{W}_{dr} \in \{0,1\}^{n \times m}$ are integrated into a heterogeneous
adjacency matrix:
%
\begin{equation}
\mathbf{T} = \begin{bmatrix} \mathbf{S}_{rr} & \mathbf{W}_{dr}^\top \\
\mathbf{W}_{dr} & \mathbf{S}_{dd} \end{bmatrix}
\in \mathbb{R}^{(m+n) \times (m+n)}
\label{eq:T}
\end{equation}

BNNR recovers missing entries via nuclear norm minimization
\cite{candes2009}:
%
\begin{equation}
\min_{\mathbf{X}} \|\mathbf{X}\|_* + \frac{\alpha}{2}
\|P_\Omega(\mathbf{W} - \mathbf{T})\|_F^2 \quad
\text{s.t.} \quad \mathbf{X} = \mathbf{W}, \quad 0 \leq \mathbf{W}_{ij} \leq 1
\label{eq:bnnr}
\end{equation}
%
where $\alpha > 0$ controls the trade-off between nuclear norm regularization
and data fidelity, $P_\Omega$ is the projection onto observed entries, and the
range constraint $[0, 1]$ ensures predicted values are interpretable as
association probabilities. BNNR solves this via ADMM \cite{boyd2011}, forming
the augmented Lagrangian:
%
\begin{equation}
\mathcal{L}_\beta(\mathbf{W}, \mathbf{X}, \mathbf{Y}) =
\|\mathbf{X}\|_* + \frac{\alpha}{2}\|P_\Omega(\mathbf{W} - \mathbf{T})\|_F^2
+ \langle \mathbf{Y}, \mathbf{W} - \mathbf{X} \rangle
+ \frac{\beta}{2}\|\mathbf{W} - \mathbf{X}\|_F^2
\label{eq:lagrangian}
\end{equation}
%
where $\beta > 0$ is the ADMM penalty parameter and $\mathbf{Y}$ is the
Lagrange multiplier. The ADMM iterations alternate between: (1) closed-form
$\mathbf{W}$-update, (2) Singular Value Thresholding (SVT; \cite{cai2010})
for the $\mathbf{X}$-update, and (3) dual ascent
$\mathbf{Y} \leftarrow \mathbf{Y} + \beta(\mathbf{W} - \mathbf{X})$.

\subsubsection{Gaussian Interaction Profile (GIP) Kernel}

The GIP kernel \cite{yamanishi2008,vanlaarhoven2011} computes similarity between
entities based on their interaction profiles:
%
\begin{align}
\mathbf{G}_{drug}(i, j) &= \exp\left(-\gamma
\|\mathbf{W}_{dr}(:, i) - \mathbf{W}_{dr}(:, j)\|^2\right) \label{eq:gip_drug} \\
\mathbf{G}_{dis}(i, j) &= \exp\left(-\gamma
\|\mathbf{W}_{dr}(i, :) - \mathbf{W}_{dr}(j, :)\|^2\right) \label{eq:gip_dis}
\end{align}
%
where $\gamma$ is the kernel bandwidth. GIP is fused with raw similarities:
$\tilde{\mathbf{S}} = w_{gip} \mathbf{G} + (1 - w_{gip}) \mathbf{S}_{raw}$.
The fusion weight $w_{gip}$ balances the contribution of association-derived
structure against chemical/clinical priors.

\subsubsection{GF-BNNR (Graph-Filtered BNNR)}

GF-BNNR consists of two steps: (1) run BNNR with GIP-fused similarities to
obtain a raw completed matrix $\mathbf{M}_{raw}$, and (2) apply a bi-directional
graph low-pass filter:
%
\begin{equation}
\mathbf{M}_{filtered} = (\mathbf{I} + \alpha_f \mathbf{L}_{dis})^{-1}
\cdot \mathbf{M}_{raw} \cdot
(\mathbf{I} + \alpha_f \mathbf{L}_{drug})^{-1}
\label{eq:gffilter}
\end{equation}
%
where $\mathbf{L} = \mathbf{I} - \mathbf{D}^{-1/2} \mathbf{S} \mathbf{D}^{-1/2}$
is the symmetric normalized graph Laplacian \cite{chung1997},
$\mathbf{D} = \text{diag}(\mathbf{S}\mathbf{1})$ is the degree matrix, and
$\alpha_f \geq 0$ controls the filter strength. The filter enforces smoothness:
the filtered prediction for a disease-drug pair is a weighted average of
predictions for similar diseases and drugs. GF-BNNR reduces to BNNR when
$\alpha_f = 0$.

\subsection{BADGE: Joint Matrix-Kernel Estimation}

\subsubsection{Motivation}

In BNNR and its \texttt{Demo2} extension, the similarity kernel $\mathbf{S}$
is a function of the \textit{sparse observation} $\mathbf{W}_{dr}$:
$\mathbf{S} = \mathbf{S}(\mathbf{W}_{dr})$. This kernel is computed once and
treated as an exogenous fixed input. However, after completing the matrix we
obtain $\mathbf{M} \succ \mathbf{W}_{dr}$ (in the sense that $\mathbf{M}$
contains richer association information), and the kernel
$\mathbf{S}(\mathbf{M})$ is strictly more informative than
$\mathbf{S}(\mathbf{W}_{dr})$. This creates an interdependence: $\mathbf{M}$
depends on $\mathbf{S}$ through the BNNR optimization, and $\mathbf{S}$ depends
on $\mathbf{M}$ through the GIP computation.

Rather than treating $\mathbf{S}$ as fixed, we formulate a joint estimation
problem over both variables (Eqs.~\ref{eq:joint_obj}--\ref{eq:kernel_constraint}).
The objective~(\ref{eq:joint_obj}) augments the standard BNNR nuclear norm
minimization with a graph regularization term that penalizes non-smoothness of
$\mathbf{M}$ with respect to the kernel $\mathbf{S}$. The
constraint~(\ref{eq:kernel_constraint}) ties the kernel to the completion:
$\mathbf{S}$ must equal the GIP kernel of the completed matrix $\mathbf{M}$,
fused with raw similarities. When $w_{gip}=0$, the kernel reduces to raw
similarities; when $w_{gip}=1$, it is pure GIP from $\mathbf{M}$.

\subsubsection{Alternating Minimization}

Problem~(\ref{eq:joint_obj})--(\ref{eq:kernel_constraint}) is biconvex:
for fixed $\mathbf{S}$, the $\mathbf{M}$-subproblem is a convex nuclear norm
minimization with graph regularization; for fixed $\mathbf{M}$, the
$\mathbf{S}$-update reduces to the closed-form GIP computation~(\ref{eq:kernel_constraint}).
We solve it via alternating minimization, initializing $\mathbf{S}^{(0)}$ from
the sparse observation and iterating:
%
\begin{align}
\mathbf{M}^{(t)} &= \arg\min_{\mathbf{M}} \;
\|\mathbf{M}\|_* + \frac{\alpha}{2}\|P_\Omega(\mathbf{M} - \mathbf{Y})\|_F^2
+ \alpha_f \cdot
\big(\text{tr}(\mathbf{M}^\top \mathbf{L}_t^{(t-1)} \mathbf{M})
+ \text{tr}(\mathbf{M} \mathbf{L}_d^{(t-1)} \mathbf{M}^\top)\big)
\label{eq:m_step} \\
\mathbf{S}_d^{(t)} &= w_{gip} \cdot \text{GIP}_{\text{drug}}(\mathbf{M}^{(t)})
+ (1 - w_{gip}) \cdot \mathbf{S}_{rr} \label{eq:s_step}
\end{align}
%
where $\mathbf{L}_d^{(t-1)}, \mathbf{L}_t^{(t-1)}$ are the normalized Laplacians
of $\mathbf{S}_d^{(t-1)}, \mathbf{S}_t^{(t-1)}$. The $\mathbf{M}$-subproblem is
solved via ADMM (Algorithm~1 in \cite{yang2019}) augmented with the graph
regularization term, and the graph filter
$(\mathbf{I} + \alpha_f \mathbf{L})^{-1}$ emerges naturally from the gradient of
the graph regularization term. The alternating scheme is summarized in
Algorithm~\ref{alg:badge}.

\subsubsection{Convergence Properties}

The objective~(\ref{eq:joint_obj}) is biconvex and the feasible set defined by~(\ref{eq:kernel_constraint})
is closed. Tseng~\cite{tseng2001} establishes that alternating minimization for
biconvex problems converges to a stationary point when each subproblem is solved
exactly and the coupling is through a continuous mapping. Both conditions hold:
the $\mathbf{M}$-subproblem~(\ref{eq:m_step}) is solved to convergence by ADMM,
and the GIP kernel map $\mathbf{M} \mapsto \text{GIP}(\mathbf{M})$ is a
composition of the Gaussian RBF kernel (Lipschitz continuous) and matrix
operations, hence continuous.

The rapid convergence ($N = 2$ suffices; Section~\ref{sec:convergence}) follows
from the structure of the GIP kernel. The kernel
$\mathbf{S} = \text{GIP}(\mathbf{M})$ depends on $\mathbf{M}$ through pairwise
Euclidean distances between association profiles. After one round of BNNR
completion and graph filtering, the resulting $\mathbf{M}^{(1)}$ is
substantially smoother than the sparse $\mathbf{W}_{dr}$, but the \textit{relative
ordering} of profile distances---which determines the GIP kernel---is already
well captured. A second kernel computation refines the absolute distance scale
but changes the kernel structure only marginally, so
$\mathbf{S}^{(2)} \approx \mathbf{S}^{(1)}$ and further iterations provide no
additional benefit.

\subsubsection{Algorithm}

\begin{algorithm}[ht]
\caption{BADGE: Bi-iterative Adaptive Drug-disease Graph Enhancement}
\label{alg:badge}
\begin{algorithmic}[1]
\Require $\mathbf{S}_{rr}$, $\mathbf{S}_{dd}$, $\mathbf{W}_{dr}$,
         $\alpha$, $\beta$, $\alpha_f$, $\gamma_{gip}$, $w_{gip}$, $N$
\State Compute GIP from sparse $\mathbf{W}_{dr}$:
        $\mathbf{G}_{drug}$, $\mathbf{G}_{dis}$
\State Fuse with raw similarities:
        \begin{align*}
        \tilde{\mathbf{S}}_{drug} &= w_{gip}\mathbf{G}_{drug} + (1-w_{gip})\mathbf{S}_{rr} \\
        \tilde{\mathbf{S}}_{dis}  &= w_{gip}\mathbf{G}_{dis}  + (1-w_{gip})\mathbf{S}_{dd}
        \end{align*}
\State $\mathbf{M}_{cur} \leftarrow \mathbf{W}_{dr}$
\For{$t = 1$ to $N$}
    \State Build $\mathbf{T}$ from
           $(\tilde{\mathbf{S}}_{drug}, \tilde{\mathbf{S}}_{dis}, \mathbf{M}_{cur})$
    \State $\mathbf{M}_{raw} \leftarrow$ BNNR($\alpha$, $\beta$, $\mathbf{T}$, $\Omega$)
    \State Compute $\mathbf{L}_{drug}$, $\mathbf{L}_{dis}$ from
           $\tilde{\mathbf{S}}_{drug}$, $\tilde{\mathbf{S}}_{dis}$
    \State $\mathbf{M}_{filtered} \leftarrow
           (\mathbf{I} + \alpha_f \mathbf{L}_{dis})^{-1} \mathbf{M}_{raw}
           (\mathbf{I} + \alpha_f \mathbf{L}_{drug})^{-1}$
    \State Preserve known entries:
           $\mathbf{M}_{cur}[known] = \mathbf{W}_{dr}[known]$
    \If{$t < N$}
        \State Recompute GIP from $\mathbf{M}_{cur}$:
                $\mathbf{G}_{drug}$, $\mathbf{G}_{dis}$
        \State Update similarities:
                $\tilde{\mathbf{S}}_{drug} = w_{gip}\mathbf{G}_{drug}
                 + (1-w_{gip})\mathbf{S}_{rr}$
        \State $\tilde{\mathbf{S}}_{dis} = w_{gip}\mathbf{G}_{dis}
                 + (1-w_{gip})\mathbf{S}_{dd}$
    \EndIf
\EndFor
\State \Return $\mathbf{M}_{cur}$
\end{algorithmic}
\end{algorithm}

BADGE with $N = 1$ is exactly GF-BNNR: the GIP recomputation branch is never
entered. With $N = 2$, BADGE performs one round of GIP refinement. Each
subsequent iteration recomputes GIP from the filtered completion of the
previous round, enabling the similarity graphs to adapt as the completion
improves. On iterations $t \geq 2$, the augmented matrix $\mathbf{T}$ is built
from $\mathbf{M}_{cur}$ (containing filtered estimates for unknown entries)
rather than the original sparse $\mathbf{W}_{dr}$, implementing a feedback loop
where improved completions inform subsequent rounds.

\subsubsection{Design Choices}

\textbf{GIP fusion weight $w_{gip}=0.3$.} In the joint formulation, $w_{gip}$
controls how strongly the kernel $\mathbf{S}$ is tied to the completion
$\mathbf{M}$ versus anchored to raw similarities. A systematic sweep across
$w_{gip} \in \{0.0, 0.1, 0.2, 0.3, 0.5, 0.7, 1.0\}$ on all three datasets
(Table~\ref{tab:wgip_sweep}) shows that $w_{gip}=0.3$ achieves near-optimal
AUPR on all datasets. Setting $w_{gip}=0$ decouples the kernel from the
completion entirely (kernel reduces to raw similarities), causing AUPR to
collapse from 0.32 to 0.13 on DNdataset. Setting $w_{gip}=1.0$ removes the raw
similarity anchor, and the kernel becomes pure GIP from the (initially sparse)
completion, collapsing to 0.03. The intermediate value $w_{gip}=0.3$ balances
the endogenous GIP signal against the exogenous raw similarity prior.

\textbf{Number of iterations $N=2$.} The alternating minimization converges at
$N=2$ across all datasets (Section~\ref{sec:convergence}), consistent with the
theoretical analysis above: the GIP kernel map
$\mathbf{M} \mapsto \text{GIP}(\mathbf{M})$ is near its fixed point after one
update. A third iteration recomputes GIP from a completion that is nearly
kernel-stationary, yielding negligible change. We recommend $N = 2$ as the
practical default.

\textbf{Relationship to prior methods.} Setting $N=1$ and $\alpha_f=0$ recovers
BNNR \texttt{Demo2} (one-shot GIP, no filter). Setting $N=1$ with $\alpha_f>0$
recovers GF-BNNR (one-shot GIP with graph filter). BADGE with $N \geq 2$ and
$\alpha_f>0$ is the full joint estimation framework with iterative kernel
refinement. This unified view shows that prior GIP-enhanced BNNR variants are
special cases of the joint formulation with early stopping at $N=1$.

\subsection{Evaluation Protocol}

We adopt 10-fold cross-validation under the CVa scheme (CV on association pairs;
random split of known drug-disease pairs into 10 equal-sized folds, with all
unknown entries treated as candidate negatives). This setting is the standard
evaluation framework in drug repositioning \cite{yang2019,luo2018}. Performance
is assessed using: (1) Area Under the ROC Curve (AUROC), (2) Area Under the
Precision-Recall Curve (AUPR), and (3) Precision@K and Recall@K for
$K \in \{10, 20\}$. Given extreme class imbalance (positive rate $\sim$1\%),
AUPR is the primary evaluation metric \cite{saito2015}. All experiments use a
fixed random seed (12345) for reproducibility.

BNNR hyperparameters follow \cite{yang2019}: $\alpha = 1$, $\beta = 10$,
convergence tolerances $tol_1 = 2 \times 10^{-3}$ and $tol_2 = 10^{-5}$,
maximum 300 iterations. BADGE-specific parameters: $\alpha_f = 0.5$,
$w_{gip} = 0.3$, $\gamma_{gip} = 1.0$, $N = 2$.

% ===========================================================================
\section{Results}
% ===========================================================================

\subsection{Overall Performance}

Table \ref{tab:main} presents the main results comparing BNNR, GBNNR, GF-BNNR,
and BADGE ($N = 2, 3$) across all three benchmark datasets.

\begin{table}[ht]
\centering
\caption{Performance comparison (10-fold CVa, mean $\pm$ std).
         Best AUPR per dataset in \textbf{bold}.}
\label{tab:main}
\small
\begin{tabular}{llcccc}
\toprule
Dataset & Method & AUROC & AUPR & P@10 & P@20 \\
\midrule
\multirow{5}{*}{\textbf{Fdataset}}
& BNNR          & 0.9319 $\pm$ 0.0157 & 0.3061 $\pm$ 0.0240 & 0.93 $\pm$ 0.05 & 0.905 $\pm$ 0.076 \\
& GBNNR         & 0.9282 $\pm$ 0.0162 & 0.3199 $\pm$ 0.0273 & 1.00 $\pm$ 0.00 & 0.975 $\pm$ 0.042 \\
& GF-BNNR       & \textbf{0.9370} $\pm$ 0.0146 & 0.3153 $\pm$ 0.0251 & 0.97 $\pm$ 0.05 & 0.925 $\pm$ 0.059 \\
& BADGE ($N{=}2$) & 0.9339 $\pm$ 0.0155 & \textbf{0.3233} $\pm$ 0.0280 & 0.99 $\pm$ 0.03 & 0.930 $\pm$ 0.054 \\
& BADGE ($N{=}3$) & 0.9299 $\pm$ 0.0165 & 0.3199 $\pm$ 0.0287 & 0.97 $\pm$ 0.05 & 0.915 $\pm$ 0.058 \\
\midrule
\multirow{5}{*}{\textbf{Cdataset}}
& BNNR          & 0.9502 $\pm$ 0.0116 & 0.2772 $\pm$ 0.1212 & 0.90 $\pm$ 0.05 & 0.945 $\pm$ 0.028 \\
& GBNNR         & 0.9472 $\pm$ 0.0118 & 0.4006 $\pm$ 0.0198 & 1.00 $\pm$ 0.00 & 0.995 $\pm$ 0.016 \\
& GF-BNNR       & \textbf{0.9545} $\pm$ 0.0133 & 0.3958 $\pm$ 0.0195 & 0.95 $\pm$ 0.05 & 0.950 $\pm$ 0.024 \\
& BADGE ($N{=}2$) & 0.9517 $\pm$ 0.0140 & \textbf{0.4051} $\pm$ 0.0215 & \textbf{1.00} $\pm$ 0.00 & \textbf{1.00} $\pm$ 0.00 \\
& BADGE ($N{=}3$) & 0.9478 $\pm$ 0.0144 & 0.4037 $\pm$ 0.0227 & 1.00 $\pm$ 0.00 & 1.00 $\pm$ 0.00 \\
\midrule
\multirow{5}{*}{\textbf{DNdataset}}
& BNNR          & 0.9288 $\pm$ 0.0172 & 0.2564 $\pm$ 0.1345 & 0.36 $\pm$ 0.14 & 0.275 $\pm$ 0.089 \\
& GBNNR         & 0.9206 $\pm$ 0.0172 & 0.2539 $\pm$ 0.1331 & 0.31 $\pm$ 0.14 & 0.220 $\pm$ 0.095 \\
& GF-BNNR       & \textbf{0.9725} $\pm$ 0.0100 & 0.3166 $\pm$ 0.0207 & 0.48 $\pm$ 0.21 & 0.345 $\pm$ 0.096 \\
& BADGE ($N{=}2$) & 0.9683 $\pm$ 0.0108 & \textbf{0.3207} $\pm$ 0.0227 & \textbf{0.50} $\pm$ 0.21 & \textbf{0.355} $\pm$ 0.083 \\
& BADGE ($N{=}3$) & 0.9638 $\pm$ 0.0120 & 0.3211 $\pm$ 0.0219 & 0.51 $\pm$ 0.20 & 0.360 $\pm$ 0.105 \\
\bottomrule
\end{tabular}
\end{table}

\textbf{BADGE ($N{=}2$) achieves the highest AUPR on all three datasets.}
Starting from the BNNR baseline, each successive method adds capability:
GBNNR incorporates manifold structure into the optimization objective
($+4.5\%$ AUPR on Fdataset, $+44.5\%$ on Cdataset); GF-BNNR applies
post-hoc graph smoothing ($+3.0\%$ on Fdataset, $+42.8\%$ on Cdataset, $+23.5\%$
on DNdataset vs.\ BNNR); and BADGE adds iterative GIP refinement, achieving
the highest AUPR across all three datasets ($+5.6\%$, $+46.1\%$, $+25.1\%$
vs.\ BNNR on Fdataset, Cdataset, and DNdataset respectively). The progressive
improvement validates the iterative approach: completing the matrix reveals
richer similarity structure, and feeding that structure back into a second
round of completion yields additional gains beyond a single pass.

On top-K precision, BADGE ($N{=}2$) achieves perfect P@10 and P@20 on Cdataset
(1.00 $\pm$ 0.00), and improves P@10 on DNdataset from 0.48 (GF-BNNR) to 0.50.

\textbf{BADGE converges at $N = 2$.} On Fdataset and Cdataset, $N{=}3$
performs slightly worse than $N{=}2$ (Fdataset: 0.3199 vs.\ 0.3233; Cdataset:
0.4037 vs.\ 0.4051), suggesting that one round of GIP refinement suffices and
further iterations may begin to overfit. On DNdataset, $N{=}3$ achieves
marginally higher AUPR (0.3211 vs.\ 0.3207), within one standard deviation. We
recommend $N = 2$ as a practical default.

\subsection{Ablation Study}
\label{sec:ablation}

To quantify the contribution of each component, we conduct two ablation
experiments across all datasets, each removing exactly one mechanism:

\begin{itemize}
    \item \textbf{A1 (No GIP):} $w_{gip} = 0$ --- raw chemical and clinical
    similarities only, without GIP fusion.
    \item \textbf{A2 (No Graph Filter):} $\alpha_f = 0$ --- BADGE without the
    Laplacian graph filter.
\end{itemize}

Table \ref{tab:ablation} presents the results.

\begin{table}[ht]
\centering
\caption{Ablation results (10-fold CVa, mean $\pm$ std).
         Full BADGE ($N{=}2$) shown for reference.}
\label{tab:ablation}
\small
\begin{tabular}{llcccc}
\toprule
Dataset & Variant & AUROC & AUPR & P@10 & P@20 \\
\midrule
\multirow{3}{*}{\textbf{Fdataset} (1.04\%)}
& BADGE (full)      & 0.9339 $\pm$ 0.0155 & \textbf{0.3233} $\pm$ 0.0280 & 0.99 $\pm$ 0.03 & 0.930 $\pm$ 0.054 \\
& A1 (no GIP)       & 0.9267 $\pm$ 0.0187 & 0.3067 $\pm$ 0.0279 & 0.97 $\pm$ 0.05 & 0.910 $\pm$ 0.057 \\
& A2 (no Filter)    & 0.9361 $\pm$ 0.0133 & 0.3144 $\pm$ 0.0259 & 0.99 $\pm$ 0.03 & 0.940 $\pm$ 0.052 \\
\midrule
\multirow{3}{*}{\textbf{Cdataset} (0.93\%)}
& BADGE (full)      & 0.9517 $\pm$ 0.0140 & \textbf{0.4051} $\pm$ 0.0215 & 1.00 $\pm$ 0.00 & 1.00 $\pm$ 0.00 \\
& A1 (no GIP)       & 0.9456 $\pm$ 0.0137 & 0.3929 $\pm$ 0.0223 & 1.00 $\pm$ 0.00 & 0.975 $\pm$ 0.042 \\
& A2 (no Filter)    & 0.9526 $\pm$ 0.0122 & 0.3729 $\pm$ 0.0718 & 0.91 $\pm$ 0.07 & 0.945 $\pm$ 0.028 \\
\midrule
\multirow{3}{*}{\textbf{DNdataset} (0.015\%)}
& BADGE (full)      & 0.9683 $\pm$ 0.0108 & \textbf{0.3207} $\pm$ 0.0227 & 0.50 $\pm$ 0.21 & 0.355 $\pm$ 0.083 \\
& A1 (no GIP)       & 0.9725 $\pm$ 0.0094 & 0.1956 $\pm$ 0.1666 & 0.38 $\pm$ 0.19 & 0.330 $\pm$ 0.126 \\
& A2 (no Filter)    & 0.9260 $\pm$ 0.0254 & 0.2844 $\pm$ 0.1004 & 0.51 $\pm$ 0.16 & 0.360 $\pm$ 0.099 \\
\bottomrule
\end{tabular}
\end{table}

Two findings emerge:

\textbf{GIP fusion is essential for prediction stability on sparse data.}
Removing GIP (A1) produces the single largest AUPR drop on DNdataset ($-0.125$,
from 0.3207 to 0.1956), with a 7.3$\times$ increase in cross-fold standard
deviation (0.1666 vs.\ 0.0227). The coefficient of variation rises from 7.1\%
to 85.1\%. This reveals that GIP-based similarity fusion is the primary
stabilizing mechanism: without it, predictions become highly unreliable across
different train-test splits in the sparse-data regime. On moderate-density
datasets, A1's AUPR degradation is smaller but consistent
($\Delta\text{AUPR} = -0.017$ and $-0.012$).

\textbf{The graph filter provides accuracy and regularization.} Removing the
filter (A2) causes the largest AUROC degradation on DNdataset
($0.9683 \to 0.9260$, $-4.4\%$), validating that Laplacian smoothing is the
principal source of predictive accuracy in the sparse regime. On Cdataset, A2
exhibits a 3.3$\times$ increase in AUPR standard deviation (0.0718 vs.\ 0.0215),
indicating that the filter provides critical regularization even on
moderate-density data.

\subsection{GIP Fusion Weight Sensitivity}

To validate the choice of $w_{gip}=0.3$, we sweep
$w_{gip} \in \{0.0, 0.1, 0.2, 0.3, 0.5, 0.7, 1.0\}$ across all three datasets
using 5-fold cross-validation.

\begin{table}[ht]
\centering
\caption{GIP fusion weight sensitivity (5-fold CVa, mean $\pm$ std).
         Default $w_{gip}=0.3$ in \textbf{bold}.}
\label{tab:wgip_sweep}
\small
\begin{tabular}{lccccc}
\toprule
$w_{gip}$ & \multicolumn{2}{c}{Fdataset} & \multicolumn{2}{c}{Cdataset} & DNdataset \\
           & AUROC & AUPR & AUROC & AUPR & AUPR \\
\midrule
0.0        & 0.9369 $\pm$ 0.0148 & 0.3091 $\pm$ 0.0295 & 0.9518 $\pm$ 0.0147 & 0.3933 $\pm$ 0.0158 & 0.1294 $\pm$ 0.1737 \\
0.1        & 0.9390 $\pm$ 0.0142 & 0.3163 $\pm$ 0.0309 & 0.9539 $\pm$ 0.0149 & 0.3985 $\pm$ 0.0148 & 0.3255 $\pm$ 0.0256 \\
0.2        & 0.9406 $\pm$ 0.0138 & 0.3216 $\pm$ 0.0270 & 0.9555 $\pm$ 0.0148 & 0.4009 $\pm$ 0.0125 & 0.3195 $\pm$ 0.0232 \\
\textbf{0.3}& \textbf{0.9421} $\pm$ 0.0133 & \textbf{0.3258} $\pm$ 0.0294 & \textbf{0.9568} $\pm$ 0.0146 & \textbf{0.4037} $\pm$ 0.0111 & \textbf{0.3175} $\pm$ 0.0243 \\
0.5        & 0.9440 $\pm$ 0.0130 & 0.3256 $\pm$ 0.0266 & 0.9587 $\pm$ 0.0140 & 0.4047 $\pm$ 0.0117 & 0.3135 $\pm$ 0.0215 \\
0.7        & 0.9439 $\pm$ 0.0132 & 0.3216 $\pm$ 0.0237 & 0.9594 $\pm$ 0.0134 & 0.3959 $\pm$ 0.0141 & 0.3075 $\pm$ 0.0209 \\
1.0        & 0.9339 $\pm$ 0.0158 & 0.2527 $\pm$ 0.0833 & 0.9528 $\pm$ 0.0093 & 0.3515 $\pm$ 0.0103 & 0.0260 $\pm$ 0.0551 \\
\bottomrule
\end{tabular}
\end{table}

The sweep reveals a consistent inverted-U pattern. At the extremes:
$w_{gip}=0.0$ (no GIP, pure raw similarities) causes catastrophic AUPR
collapse on DNdataset ($0.32 \to 0.13$), demonstrating that GIP is
indispensable for sparse-data prediction. $w_{gip}=1.0$ (pure GIP, no raw
similarities) degrades performance on all datasets and collapses to near-zero
AUPR on DNdataset ($0.03$), showing that raw chemical and clinical
similarities provide essential ground truth. The default $w_{gip}=0.3$ lies
near the AUPR peak on all datasets, validating this choice as a robust default.

\subsection{Cross-Dataset Variance Reduction}
\label{sec:variance}

A notable pattern in Table~\ref{tab:main} is the reduction in cross-fold AUPR
variability for methods that incorporate graph structure. On DNdataset, BNNR
exhibits AUPR standard deviation of 0.1345 (CV $\approx$ 52\%); GF-BNNR
reduces this to 0.0207 (CV $\approx$ 6.5\%), and BADGE ($N{=}2$) to 0.0227
(CV $\approx$ 7.1\%). Conversely, removing GIP (A1,
Section~\ref{sec:ablation}) inflates AUPR standard deviation to 0.1666
(CV $\approx$ 85\%), the worst cross-fold stability of any variant tested. This
confirms that GIP fusion, not merely the graph filter, is the primary source of
prediction robustness. The order-of-magnitude reduction in variability achieved
by manifold-based methods indicates that they provide substantially more
reliable predictions across different train-test splits---a practically
valuable property when deploying predictions for experimental validation.

\subsection{Convergence Analysis}
\label{sec:convergence}

Table~\ref{tab:convergence} shows BADGE's performance as a function of the
number of refinement iterations $N$.

\begin{table}[ht]
\centering
\caption{Effect of refinement iterations on AUPR (10-fold CVa).}
\label{tab:convergence}
\begin{tabular}{cccc}
\toprule
$N$             & Fdataset AUPR & Cdataset AUPR & Interpretation \\
\midrule
1 (GF-BNNR)     & 0.3153        & 0.3958        & No refinement \\
2               & 0.3233        & 0.4051        & One round of GIP refinement \\
3               & 0.3199        & 0.4037        & No further benefit \\
\bottomrule
\end{tabular}
\end{table}

BADGE converges in two iterations. The third iteration yields AUPR values
between $N{=}1$ and $N{=}2$, suggesting that the second round of GIP
computation begins to overfit to the completion. We recommend $N = 2$ as the
default. The first iteration (GF-BNNR) produces a completion already
substantially better than the sparse input. The second iteration refines GIP
using this improved completion. However, the improvement is bounded: the
filtered completion, while better than the sparse input, cannot produce a GIP
that is arbitrarily more informative than the original. A third iteration
refines using a GIP only marginally different from the second, providing no
additional signal.

\begin{figure}[ht]
\centering
\includegraphics[width=0.95\textwidth]{figures/fig1_main_results.pdf}
\caption{AUPR comparison across methods and datasets (10-fold CVa,
         mean $\pm$ standard deviation over 10 folds).
         BADGE ($N{=}2$) achieves the highest AUPR on all three datasets.}
\label{fig:main_results}
\end{figure}

\begin{figure}[ht]
\centering
\includegraphics[width=0.95\textwidth]{figures/fig2_framework.pdf}
\caption{BADGE framework schematic. Initial GIP is computed from sparse
         associations and fused with raw similarities. The iterative
         refinement loop (dashed boundary) repeats for $N$ iterations,
         recomputing GIP from the filtered completion after each round.}
\label{fig:framework}
\end{figure}

\begin{figure}[ht]
\centering
\includegraphics[width=0.95\textwidth]{figures/fig3_convergence.pdf}
\caption{Convergence analysis: AUPR as a function of refinement
         iterations $N$. BADGE converges at $N=2$ on moderate-density
         datasets.}
\label{fig:convergence}
\end{figure}

% ===========================================================================
\section{Discussion}
% ===========================================================================

\subsection{Why Joint Estimation Improves over Fixed-Kernel BNNR}

BNNR \cite{yang2019} treats the similarity kernel as exogenous: the drug-drug
and disease-disease matrices are fixed inputs, computed once from external data
(chemical fingerprints, MeSH term overlap) or, in its \texttt{Demo2} extension,
augmented with one-shot GIP from sparse associations. The optimization
objective~(\ref{eq:bnnr}) only sees the \textit{values} of the kernel, not its
\textit{provenance}. GBNNR and GF-BNNR add manifold regularization via the graph
Laplacian of this fixed kernel, improving smoothness of the solution with respect
to the given graph, but the kernel itself remains fixed.

The joint formulation~(\ref{eq:joint_obj})--(\ref{eq:kernel_constraint})
fundamentally changes the role of the kernel: it becomes an optimization
variable that adapts to the improving completion. At iteration $t=1$, the kernel
$\mathbf{S}^{(0)}$ is computed from sparse $\mathbf{W}_{dr}$---identical to
\texttt{Demo2} and GF-BNNR. At iteration $t=2$, the kernel $\mathbf{S}^{(1)}$ is
computed from $\mathbf{M}^{(1)}$, the filtered completion from iteration 1. This
kernel is strictly more informative because $\mathbf{M}^{(1)}$ contains dense
estimates for all entries, revealing similarity structure invisible in the sparse
$\mathbf{W}_{dr}$ (which is 99.985\% zeros in DNdataset). The second BNNR
completion benefits from this improved kernel, yielding the additional AUPR
gains over GF-BNNR observed in Table~\ref{tab:main}.

This also explains why the improvement is larger on sparser datasets
($+2.5\%$ on DNdataset vs.\ $+2.5\%$ on Fdataset for BADGE over GF-BNNR):
the denser the original data, the less additional information the first
completion provides for GIP recomputation.

\subsection{Why Two Iterations Suffice}

The convergence at $N=2$ follows from the fixed-point structure of the GIP
kernel update. Define the kernel update operator
$\mathcal{K}: \mathbf{M} \mapsto \mathbf{S}$ via
Eq.~(\ref{eq:kernel_constraint}), and the completion operator
$\mathcal{C}: \mathbf{S} \mapsto \mathbf{M}$ via
Eq.~(\ref{eq:m_step}). The alternating minimization iterates
$\mathbf{S}^{(t)} = \mathcal{K}(\mathcal{C}(\mathbf{S}^{(t-1)}))$, which is a
fixed-point iteration on the composite map $\mathcal{K} \circ \mathcal{C}$.

The GIP kernel~(\ref{eq:gip_drug})--(\ref{eq:gip_dis}) is a Gaussian RBF over
pairwise profile distances. The profile of entity $i$ is the $i$-th row or column
of the association matrix. After one completion, these profiles change from
binary sparse vectors to dense real-valued vectors, substantially compressing the
distance distribution. However, after the second completion, the profiles change
only incrementally---the completion has already converged to within the BNNR
tolerance---so the GIP kernel computed from $\mathbf{M}^{(2)}$ is nearly
identical to that from $\mathbf{M}^{(1)}$. Formally,
$\|\mathbf{S}^{(2)} - \mathbf{S}^{(1)}\|_F \approx 0$, and additional iterations
leave the kernel unchanged up to numerical precision.

This analysis also yields a practical diagnostic: if
$\|\mathbf{S}^{(t)} - \mathbf{S}^{(t-1)}\|_F$ falls below a threshold (e.g.,
$10^{-3}$), the alternating minimization has converged and further BNNR calls can
be skipped. In our experiments, this condition is satisfied at $t=2$ on all
datasets.

\subsection{The Role of GIP Fusion}

The $w_{gip}$ sensitivity analysis (Table~\ref{tab:wgip_sweep}) reveals that
GIP fusion is not merely a performance enhancer but an \textit{essential}
component. On DNdataset, removing GIP ($w_{gip}=0$) causes AUPR to drop from
0.32 to 0.13, with cross-fold standard deviation exploding from 0.02 to 0.17.
This demonstrates that on ultra-sparse data, raw chemical and clinical
similarities alone are insufficient---the manifold structure provided by GIP
is necessary for both accuracy and stability.

Conversely, pure GIP ($w_{gip}=1.0$) degrades performance everywhere and
collapses to 0.03 AUPR on DNdataset, confirming that GIP must be grounded by
raw similarities. The intermediate value $w_{gip}=0.3$---giving GIP a minority
but non-trivial weight---optimally balances these complementary sources.

\subsection{Limitations and Future Work}
\label{sec:limitations}

Several limitations should be noted. First, the current experiments use
different SVD strategies in the BNNR inner loop: GF-BNNR uses full SVD while
BADGE uses adaptive truncated SVD, which is approximately 30$\times$ faster on
large matrices. While both strategies produce mathematically equivalent
solutions at convergence, the adaptive variant may reach slightly different
results due to numerical tolerances. This difference should be harmonized in
future experiments.

Second, the current study lacks formal statistical significance testing
(e.g., paired Wilcoxon test or Friedman test with post-hoc Nemenyi). While the
consistent AUPR improvement across all three datasets is suggestive, the
overlapping standard deviations mean that per-dataset differences are not
individually significant at the 95\% confidence level. A multi-dataset
meta-analysis would provide stronger statistical support.

Third, BADGE currently uses a single GIP fusion weight $w_{gip}$ for both drug
and disease sides. In principle, drug-side and disease-side similarities could
have different reliability and warrant separate fusion weights. This extension
is straightforward and merits investigation.

Fourth, while $w_{gip}=0.3$ performs well across all three benchmark datasets,
its generalizability to datasets with different characteristics (e.g.,
different similarity computation methods, different association patterns) has
not been established. A more systematic calibration across a broader range of
datasets would strengthen this recommendation.

Future work should explore: (1) per-entity adaptive GIP fusion based on local
neighborhood density rather than a single global $w_{gip}$; (2) extension to
multi-relational heterogeneous networks incorporating drug-target and
protein-protein interaction data; (3) theoretical analysis of the optimal
number of refinement iterations as a function of matrix dimensions and density;
and (4) application of the iterative GIP refinement framework to other matrix
completion methods beyond BNNR.

% ===========================================================================
\section{Conclusion}
% ===========================================================================

We have presented BADGE, a joint matrix-kernel estimation framework for drug
repositioning that reformulates the role of GIP similarity from a fixed
preprocessing step to an endogenous optimization variable. The key insight is
that the completion matrix and the similarity kernel are mutually dependent, and
alternating minimization between them converges in two iterations to a
stationary point of the biconvex joint objective. BADGE achieves the highest
AUPR on all three benchmark datasets, nests BNNR \texttt{Demo2} and GF-BNNR as
special cases with $N=1$, and requires no per-dataset hyperparameter tuning
beyond the single global default $w_{gip}=0.3$.

The alternating minimization framework provides a principled alternative to the
prevalent fixed-kernel paradigm. The convergence at $N=2$ is not an empirical
coincidence but a consequence of the GIP kernel map stabilizing after one
refinement, as the filtered completion is already near the kernel fixed point.
This analysis offers both theoretical justification for the observed performance
and a practical diagnostic ($\|\mathbf{S}^{(t)} - \mathbf{S}^{(t-1)}\|_F$) for
determining convergence.

The joint estimation framework extends beyond BNNR: any matrix completion method
that uses a similarity kernel computed from the association matrix can in
principle benefit from treating that kernel as an endogenous variable.
Recomputing the kernel from the completed output and iterating is a
general-purpose strategy for strengthening the coupling between completion and
side information.

% ===========================================================================
\section*{Data Availability}
% ===========================================================================

The datasets used in this study are publicly available from prior publications
\cite{gottlieb2011,yang2019}. The source code for BADGE is available at
\url{https://github.com/yankangyk/bnnr-innovation}.

% ===========================================================================
\section*{Funding}
% ===========================================================================

This work was supported by the National Natural Science Foundation of China
[grant number pending].

% ===========================================================================
\section*{Conflict of Interest}
% ===========================================================================

None declared.

% ===========================================================================
\section*{Author Contributions}
% ===========================================================================

Y.K.\ conceived the method, implemented the code, conducted experiments, and
wrote the manuscript.

% ===========================================================================
\section*{Acknowledgments}
% ===========================================================================

The authors thank the anonymous reviewers for their constructive feedback.

% ===========================================================================
% BIBLIOGRAPHY
% ===========================================================================
\bibliographystyle{oup-abbrvnat}
\bibliography{references}

\end{document}
